\section{数据表示与运算}

\subsection{数制与编码}

计算机中所有信息都以二进制（0/1）存储。本章讨论如何用二进制表示整数、浮点数、字符。

\begin{table}[H]
\centering
\caption{常用数制对比}
\label{tab:number_systems}
\begin{tabulary}{\textwidth}{CCC}
\toprule
\textbf{进制} & \textbf{基数} & \textbf{数字符号} \\
\midrule
二进制（Binary） & 2 & 0, 1 \\
八进制（Octal） & 8 & 0–7 \\
十进制（Decimal） & 10 & 0–9 \\
十六进制（Hex） & 16 & 0–9, A–F \\
\bottomrule
\end{tabulary}
\end{table}

\begin{equation}
(N)_r = \sum_{i=-m}^{n-1} d_i \cdot r^i
\end{equation}

\subsection{整数表示}

\subsubsection{无符号整数}

\begin{equation}
U(X) = \sum_{i=0}^{n-1} x_i \cdot 2^i, \quad \text{范围: } [0, 2^n - 1]
\end{equation}

\subsubsection{原码（Sign-Magnitude）}

最高位为符号位（0 正 / 1 负），其余位表示绝对值。存在 $+0$（0000）和 $-0$（1000）两个零。

\subsubsection{补码（Two's Complement）}

现代计算机统一采用补码表示有符号整数——优势是加减法统一，不存在双零。

\begin{definition}[补码]
$n$ 位有符号整数 $X$ 的补码值为：
\[
T(X) = -x_{n-1} \cdot 2^{n-1} + \sum_{i=0}^{n-2} x_i \cdot 2^i
\]
范围：$[-2^{n-1}, 2^{n-1}-1]$（如 8 位为 $[-128, 127]$）。
\end{definition}

\begin{example}[补码转换]
8 位补码 \texttt{1000 0001} 的值：
\[
-1 \cdot 2^7 + 1 \cdot 2^0 = -128 + 1 = -127
\]
取反加一法（快速转换）：原码 $\to$ 逐位取反 $\to$ 末位加一。
\end{example}

\subsubsection{移码（Biased Notation）}

用于浮点数阶码。$E_{\text{biased}} = E + \text{Bias}$，单精度 Bias = 127，双精度 Bias = 1023。

\subsection{整数运算}

\subsubsection{加法器}

\[
\text{全加器: } S = A \oplus B \oplus C_{\text{in}}, \quad C_{\text{out}} = AB + (A \oplus B)C_{\text{in}}
\]

行波进位加法器（Ripple Carry Adder）：$n$ 个全加器串联，延迟 $O(n)$。超前进位加法器（Carry Lookahead Adder）：并行计算进位，延迟 $O(\log n)$。

\subsubsection{溢出检测}

\begin{itemize}
    \item 同号相加得异号 $\to$ 溢出
    \item 补码：$A_{n-1}=B_{n-1}=0$ 且 $S_{n-1}=1$；或 $A_{n-1}=B_{n-1}=1$ 且 $S_{n-1}=0$
    \item 硬件判断：$O = C_n \oplus C_{n-1}$
\end{itemize}

\subsection{浮点数表示：IEEE 754}

\begin{definition}[IEEE 754 单精度浮点数]
32 位：1 位符号 $S$ + 8 位移码阶码 $E$ + 23 位尾数 $M$（隐含 1）：
\[
(-1)^S \times (1.M) \times 2^{E - 127}
\]
\end{definition}

\begin{table}[H]
\centering
\caption{IEEE 754 格式参数}
\label{tab:ieee754}
\begin{tabulary}{\textwidth}{CCCC}
\toprule
\textbf{精度} & \textbf{总位数} & \textbf{阶码位} & \textbf{尾数位} \\
\midrule
单精度（float） & 32 & 8 & 23 \\
双精度（double） & 64 & 11 & 52 \\
扩展双精度 & 80 & 15 & 64 \\
\bottomrule
\end{tabulary}
\end{table}

\begin{itemize}
    \item 规格化数：$E \neq 0$ 且 $E \neq 255$，尾数隐含 1
    \item 非规格化数（denormal）：$E = 0$，尾数隐含 0，用于表示极小数
    \item 无穷大：$E = 255$，$M = 0$
    \item NaN：$E = 255$，$M \neq 0$
\end{itemize}

\begin{equation}
\text{舍入模式: 向偶数舍入（Round to Nearest Even）}
\end{equation}

\subsection{浮点运算}

\begin{equation}
\begin{aligned}
\text{加法: } & \text{对阶} \rightarrow \text{尾数相加} \rightarrow \text{规格化} \rightarrow \text{舍入} \\
\text{乘法: } & \text{阶码相加} \times \text{尾数相乘} \rightarrow \text{规格化} \rightarrow \text{舍入}
\end{aligned}
\end{equation}

\subsection{字符编码}

\begin{itemize}
    \item \textbf{ASCII}：7 位（0–127），扩展 ASCII 为 8 位
    \item \textbf{Unicode}：统一的全球字符集
    \item \textbf{UTF-8}：变长编码，ASCII 兼容（1–4 字节），互联网标准
    \item \textbf{GBK / GB 2312}：中文编码标准
\end{itemize}
